\subsection{Constructing Procedures Using \code{lambda}}
\label{Section 1.3.2}

In using \code{sum} as in \link{Section 1.3.1}, it seems terribly awkward to
have to define trivial procedures such as \code{pi\-/term} and \code{pi\-/next}
just so we can use them as arguments to our higher-order procedure.  Rather
than define \code{pi\-/next} and \code{pi\-/term}, it would be more convenient to
have a way to directly specify ``the procedure that returns its input
incremented by 4'' and ``the procedure that returns the reciprocal of its input
times its input plus 2.''  We can do this by introducing the special form
\code{lambda}, which creates procedures.  Using \code{lambda} we can describe
what we want as

\begin{codeBlock}
lambda x: x + 4
\end{codeBlock}

\noindent
and

\begin{codeBlock}
lambda x: 1.0 / (x * (x + 2))
\end{codeBlock}

\noindent
Then our \code{pi\_sum} procedure can be expressed without defining any
auxiliary procedures as

\begin{codeBlock}
>>> def pi_sum(a, b):
...     return sum(lambda x: 1.0 / (x * (x + 2)),
...                a,
...                lambda x: x + 4,
...                b)
\end{codeBlock}

\noindent
Again using \code{lambda}, we can write the \code{integral} procedure without
having to define the auxiliary procedure \code{add\_dx}:

\begin{codeBlock}
>>> def integral(f, a, b, dx):
...     return sum(f,
...                a + dx / 2.0,
...                lambda x: x + dx,
...                b) * dx
\end{codeBlock}

\noindent
In general, \code{lambda} is used to create procedures in the same way as
\code{def}, except that no name is specified for the procedure:

\begin{syntaxForm}
lambda \meta{formal parameters}: \meta{body}
\end{syntaxForm}

\noindent
The resulting procedure is just as much a procedure as one that is created
using \code{def}.  The only difference is that it has not been associated
with any name in the environment.  In fact,

\begin{codeBlock}
>>> def plus4(x):
...     return x + 4
\end{codeBlock}

\noindent
is equivalent to

\begin{codeBlock}
>>> plus4 = lambda x: x + 4
\end{codeBlock}

\noindent
We can read a \code{lambda} expression as follows:

\begin{codeBlock}
lambda                       x:       x  +   4
    |                        |       |   |   |
the procedure of an argument x that adds x and 4
\end{codeBlock}

\noindent
Like any expression that has a procedure as its value, a \code{lambda}
expression can be used as the operator in a combination such as

\begin{codeBlock}
>>> (lambda x, y, z: x + y + square(z))(1, 2, 3)
~\textit{12}~
\end{codeBlock}

\noindent
or, more generally, in any context where we would normally use a procedure
name.\footnote{It would be clearer and less intimidating to people learning
Lisp if a name more obvious than \code{lambda}, such as \code{make\-/procedure},
were used.  But the convention is firmly entrenched.  The notation is adopted
from the λ-calculus, a mathematical formalism introduced by the
mathematical logician Alonzo \link{Church (1941)}.  Church developed the
λ-calculus to provide a rigorous foundation for studying the
notions of function
and function application.  The λ-calculus has become a basic tool
for mathematical investigations of the semantics of programming languages.}

\subsubsection*{Using \code{let} to create local variables}

Another use of \code{lambda} is in creating local variables.  We often need
local variables in our procedures other than those that have been bound as
formal parameters.  For example, suppose we wish to compute the function

$$ f(x,y) = x(1 + xy)^2 + y(1 - y) + (1 + xy)(1 - y), $$

\noindent
which we could also express as

$$
\begin{array}{r@{{}={}}l}
  a 	  &  1 + xy, \\
  b 	  &  1 - y,  \\
  f(x,y)  &  xa^2 + yb + ab.
\end{array}
$$

In writing a procedure to compute \( f \), we would like to include as local
variables not only \( x \) and \( y \) but also the names of intermediate
quantities like \( a \) and \( b \).  One way to accomplish this is to use an
auxiliary procedure to bind the local variables:

\begin{codeBlock}
>>> def f(x, y):
...     def f_helper(a, b):
...         return x * square(a) + y * b + a * b
...     return f_helper(1 + x * y, 1 - y)
\end{codeBlock}

\noindent
That works, but it is a roundabout way of saying something simple.  What we
want is to give two names values and then use them, and in Python we say so
directly:

\begin{codeBlock}
>>> def f(x, y):
...     a = 1 + x * y
...     b = 1 - y
...     return x * square(a) + y * b + a * b
\end{codeBlock}

\noindent
This is the version to write.  A name introduced by assignment inside a
procedure is local to that procedure: it is created when the assignment runs
and it goes away when the procedure returns, so \code{a} and \code{b} here are
as private to \code{f} as \code{x} and \code{y} are.

Lisp cannot do this, having no assignment of the kind we are using---we shall
see in \link{Chapter 3} why the original is in no hurry to introduce one---and
so it reaches instead for a special form called \code{let}:

\begin{syntaxForm}
(define (f x y)
  (let ((a (+ 1 (* x y)))
        (b (- 1 y)))
    (+ (* x (square a))
       (* y b)
       (* a b))))
\end{syntaxForm}

\noindent
which says: let \meta{var} have the value \meta{exp}, and so on for each pair,
in the body that follows.

It would be easy to leave the matter there, with two languages that spell the
same idea differently.  But the original does something more interesting with
\code{let}, and it is worth following, because it says something about
assignment that is not obvious.

\begin{buildFromScarcity}
Where does a local name come from?

The original answers by pointing out that \code{let} is not a new mechanism at
all.  It is shorthand: a \code{let} expression means exactly

\begin{syntaxForm}
(lambda \meta{var1}, ..., \meta{varn}: \meta{body})(
    \meta{exp1}, ..., \meta{expn})
\end{syntaxForm}

\noindent
a procedure of those names, applied to those values.  Nothing about local
naming is primitive.  A name is local because it is a \emph{parameter}, and a
parameter gets its value because the procedure was \emph{applied}.  Binding a
name and calling a procedure are the same act.

The \code{f} above can be written that way in Python as easily as in Lisp:

\begin{codeBlock}
>>> def f(x, y):
...     return (lambda a, b: x * square(a) + y * b + a * b)(
...         1 + x * y, 1 - y
...     )
\end{codeBlock}

\noindent
which nobody would write, and which is not being recommended.  It is worth
seeing once because it makes a claim that the assignment version conceals:
\code{a = 1 + x * y} is doing what a procedure call does, and the reason
\code{a} is local to \code{f} is the same reason \code{x} is.  When we come to
say what an environment is, in \link{Chapter 3}, and to write an interpreter
that maintains one, in \link{Chapter 4}, the frame that holds \code{a} will be
the frame made by calling a procedure---because that is what it always was.

The differences between the two forms are real, and the original sets out two
of them below.  Both are consequences of the same thing: the lambda form binds
its names for the extent of one expression, an assignment binds them for the
rest of the procedure.
\end{buildFromScarcity}

We can see from this equivalence that the scope of a name bound this way is
the body of the procedure that binds it -- one expression, and no more. This
implies two things, and in both the assignment we prefer behaves differently:

\begin{itemize}
\item
It binds a name as locally as possible to where it is used.  For example, if
the value of \code{x} is 5, the value of the expression

\begin{codeBlock}
>>> (lambda x: x + x * 10)(3) + x
~\textit{38}~
\end{codeBlock}

\noindent
is 38.  Here, the \code{x} inside the procedure is 3, so that part of the
expression is 33.  The \code{x} added to it is still 5: the binding lasted for
one expression and did not disturb what was outside it.  Writing \code{x = 3}
instead would have replaced the outer \code{x} for the rest of the procedure,
and the answer would have been 36.

\item
The values are computed outside.  This matters when an expression providing a
value depends on a name that is about to be rebound.  For example, if the value
of \code{x} is 2, the expression

\begin{codeBlock}
>>> (lambda x, y: x * y)(3, x + 2)
~\textit{12}~
\end{codeBlock}

\noindent
has the value 12 because the arguments are worked out before the procedure is
entered: \code{y} is the \emph{outer} \code{x} plus 2, which is 4, and only
then does \code{x} become 3.  Written as two assignments, \code{x = 3}
followed by \code{y = x + 2}, the second would see the new \code{x}, making
\code{y} equal to 5 and the product 15.
\end{itemize}

\noindent
Sometimes we can use internal definitions to get the same effect as with
\code{let}.  For example, we could have defined the procedure \code{f} above as

\begin{codeBlock}
>>> def f(x, y):
...     a = 1 + x * y
...     b = 1 - y
...     return x * square(a) + y * b + a * b
\end{codeBlock}

\noindent
The original prefers \code{let} in situations like this, and uses internal
\code{define} only for internal procedures.  We have no such choice to make,
and the assignment above is what we shall write.\footnote{Understanding
internal definitions well enough to be sure a program means what we intend it
to mean requires a more elaborate model of the evaluation process than we have
presented in this chapter.  The subtleties do not arise with internal
definitions of procedures, however.  We will return to this issue in
\link{Section 4.1.6}, after we learn more about evaluation.}

\begin{exerciseGroup}
\phantomsection\label{Exercise 1.34}
\exerciseHeading{Exercise 1.34:} Suppose we define the procedure

\begin{codeBlock}
>>> def f(g):
...     return g(2)
\end{codeBlock}

Then we have

\begin{codeBlock}
>>> f(square)
~\textit{4}~
>>> f(lambda z: z * (z + 1))
~\textit{6}~
\end{codeBlock}

What happens if we (perversely) ask the interpreter to evaluate
\code{f(f)}?  Explain.
\end{exerciseGroup}

